\chapter{Accuracy: The Independent Evidence}
\label{ch:accuracy}

Accuracy is a property of a specified quantity, instrument configuration and operating condition.
Ball speed at separation, club speed at a face point before impact, and predicted carry are different measurands with different validation requirements.
A study can establish useful agreement for some of these without validating every number on a display.
The evidence below separates comparisons against an independently characterized reference, repeated-swing reliability, and agreement between commercial devices.

\section{Comparison Against an Optical Reference}

Leach, Forrester, Mears and Roberts~\cite{leach2017} recorded 240 shots from eight golfers using a driver, 7-iron and utility wedge, comparing TrackMan Pro~IIIe and Foresight GC2+HMT with a four-camera 5,400\,Hz GOM system.
The experiment used the then-current software in January 2015.
The authors characterized the reference with a calibration artefact, position tests, simulated noise propagation and a separate spinning-sphere experiment.
Their simulated ball-parameter uncertainties were 0.06\,mph for speed, $0.01\degs$ for launch angle and direction, and 14\,rpm for spin; the spinning-sphere comparison gave a 9\,rpm uncertainty.
These are setup-specific estimates under the paper's stated uncertainty convention, not proof of an error-free reference.
The same propagation experiment was not performed for the club quantities.

Ball parameters generally agreed more closely than club parameters, but the details resist a simple radar-versus-camera ranking.
Table~4 of the paper gives pooled club-speed median differences (device minus reference) of $-1.1$\,mph for TrackMan and $+2.8$\,mph for Foresight, with interquartile intervals of $[-1.9,0.1]$ and $[1.4,4.5]$\,mph.
TrackMan's often-quoted $-0.4$\,mph club-speed median is the \emph{driver-only} result in Table~5.
Its $0.0\degs$ launch-angle median is also driver-only; the pooled median is $0.1\degs$.
Pooling or switching clubs changes the question being answered.

The paper's ``research grade'' fractions also need their original definition.
The tolerance interval was centered on each sample's \emph{median difference}, rather than on zero difference from the reference.
Thus the reported 54\% of TrackMan and 29\% of Foresight club-speed values within the $\pm1$\,mph research band describe spread after centering.
They are not the fractions of all attempted shots with absolute club-speed error below 1\,mph.
The thresholds were selected by the authors for research and coaching purposes; they are not a universal instrument standard.
Report bias, spread, tolerance definition and denominator together.

Availability was a separate limitation.
TrackMan returned ball and club-speed data for about 98\% of trials, but the other club parameters for 62\% overall and only 19\% of utility-wedge trials.
Foresight tracked the ball in 90\% and the club in 75\% of trials.
The simultaneous laboratory arrangement required compromises in lighting, markers and nearby equipment that the authors identified as possible interference sources.
These return rates characterize that experiment, not every current model or installation.
TrackMan spin values flagged as calculated were excluded from the analysis.

Several disagreements also concern definitions.
The reference club velocity used a manufacturer-supplied centre-of-gravity location, while a launch monitor may track a different point on a rotating head.
The study compared different face-normal definitions for the two devices and could not exactly reproduce all vendors' impact-time conventions.
Such differences are physically consequential, not merely interchangeable labels.
In this experiment, some impact-model-derived orientation results agreed more closely than optical orientation results.
The lesson is to validate each estimator and measurand; ``measured'' does not automatically mean accurate and ``derived'' does not automatically mean poor.

\section{Reliability Across Repeated Swings}

Bishop et al.~\cite{tm4reliability}, in a 2023 journal volume published online in 2024, studied 13 male golfers performing ten shots with their own driver in each of two indoor sessions.
The paper reported high within-session intraclass correlation coefficients (ICCs) for club speed (0.99) and ball speed (0.97--0.99), while spin ICCs were 0.60 and 0.02; the between-session spin ICC was 0.15.
These findings concern repeatability of the combined golfer, protocol and instrument.
Without a simultaneous independent spin reference, they do not isolate radar spin error from actual variation in impact and ball launch.
Nor does a high speed ICC establish a small absolute speed bias.

An ICC compares variance components and depends on the variation among the subjects or shots being compared.
In a simple random-intercept model, it has the form $\sigma_b^2/(\sigma_b^2+\sigma_w^2)$, where $\sigma_b^2$ is between-subject variance and $\sigma_w^2$ includes within-subject variation and measurement error.
This schematic formula is not a substitute for specifying the fitted ICC model, agreement definition and single- or average-measurement target.
A more heterogeneous group can produce a higher ICC without any improvement in the instrument.
For monitoring one golfer, absolute uncertainty and the repeatability of the chosen session summary matter alongside relative reliability.

The Bishop tables label spin dispersion in degrees per second and the abstract mislabels carry dispersion in mph, while its tables use yards for carry.
Do not silently reinterpret the spin units as rpm or propagate the carry typo into a device specification.
The source supports the reported ICC comparison, but the ambiguous spin units limit reuse of its dimensional error figures.

The uncertainty of a change also involves both measurements.
For two session estimates with error standard deviations $s_1,s_2$ and error covariance $c_{12}$,
\begin{equation}
\label{eq:changeuncertainty}
 s_{\Delta}^2=s_1^2+s_2^2-2c_{12}.
\end{equation}
Under an independent, equal-precision, approximately normal error model with known error standard deviation $s$, a two-sided 95\% noise threshold is about $1.96\sqrt{2}\,s$, not $2s$.
Estimating that standard deviation from limited data introduces additional uncertainty and requires a corresponding interval method.
This derivation assumes that $s$ describes the actual session estimator being compared and excludes systematic drift.
It does not make every larger change practically important or prove that training caused the change.
Repeated-swings variability, alignment changes and instrument error all need consideration when using such a threshold.

\section{Agreement Between Commercial Monitors}

Bliss and Langdown's full paper~\cite{mevoplusstudy} is available, replacing the earlier abstract-only account.
Published online in December 2025 in a 2026 volume, it compares Mevo+ with TrackMan~4 indoors using one golfer and three clubs.
Of 118 driver, 118 seven-iron and 174 pitching-wedge attempts, 58, 80 and 69 were retained after exclusions for missing fields, estimated spin or severe mishits: approximately 49\%, 68\% and 40\%.
These are retained analysis fractions, not either device's standalone detection rate.
Table~2 reports a seven-iron ball-speed consistency ICC of 0.996 alongside a mean TrackMan-minus-Mevo+ difference of $-3.03$\,mph; its absolute-agreement ICC is 0.90 with a wide interval from $-0.02$ to 0.97.
The protocol used mean-rating ICCs and reported both consistency and absolute agreement.
Results describe the selected paired observations from that golfer and setup; they do not establish error against an independent physical reference or performance on the excluded attempts.
The paper's criterion of 95\% of differences inside computed 95\% agreement limits does not specify a practically acceptable error width.

More generally, consistency can remain high despite an additive calibration offset.
An average-measures ICC targets an average of ratings, rather than the reliability of one instrument reading; its value should not be quoted as a single-reading accuracy specification.
Limits of agreement and confidence intervals answer different questions: the former describe the spread of paired differences, while the latter quantify uncertainty in estimated statistics.
For approximately normal differences, $\bar d\pm1.96s_d$ estimates population agreement limits; their width must be compared with a tolerance chosen for the intended use.
Merely finding most observations inside those fitted limits cannot establish interchangeability.

If two devices measure the same scalar measurand with errors $e_A$ and $e_B$, the variance of their difference is
\begin{equation}
\label{eq:paireddevicevariance}
\operatorname{Var}(e_A-e_B)
=\operatorname{Var}(e_A)+\operatorname{Var}(e_B)
 -2\operatorname{Cov}(e_A,e_B).
\end{equation}
Agreement alone cannot identify both devices' error variances or expose a bias they share.
This reasoning presumes matched reference points, timing and definitions; otherwise the physical measurand difference contributes too.
Comparing two indoor carry predictions adds another layer because both outputs depend on launch estimates and flight models.

\section{Robot Tests, Specifications and Causal Interpretation}

Robot and industry tests can be useful when the protocol and raw paired observations are available.
A smaller standard deviation of displayed spin over robot shots is not, by itself, evidence of smaller spin-measurement error.
True impact variation, ball differences, rejection rules, averaging and output smoothing can all change that standard deviation.
A constant but biased output would have excellent repeatability and poor accuracy.
Comparing these output standard deviations cannot establish a universal photometric advantage without an independent reference and an uncertainty analysis of the test.
The industry reports cited elsewhere in this review~\cite{golflabsrobot,mygolfspyindoor} should be read as protocol-specific comparisons, not replacements for that analysis.
Likewise, a manufacturer's tolerance requires its conditions, confidence convention and parameter definition before it can be compared with a study's spread statistic.

\section{A Validation Protocol That Supports the Intended Claim}

\begin{enumerate}
\item \textbf{Define the measurand.}
Specify units, coordinate frame, reference point, event time and whether carry is observed or simulated.
For club speed, a constant per-club offset is not a substitute for transforming velocities between points on a rotating rigid body (\cref{sec:pathreference,app:screw}).
\item \textbf{Characterize the reference.}
Calibrate space and time and document reference uncertainty over the tested range.
A consumer monitor can be a useful comparator, but agreement with it does not certify physical accuracy.
RPT optical markings and RCT radar-reflective construction serve different measurement routes; choose balls supported by each instrument's actual spin mode (\cref{app:impl}).
\item \textbf{Retain every attempt.}
Report attempted shots, returned values, rejected values and estimated fallbacks separately for each parameter and condition.
Analyze conditional accuracy among returned observations alongside total availability.
Missingness may depend on club, speed, strike location and spin; neither dropping failures nor assigning them zero error is defensible.
\item \textbf{Separate calibration from validation.}
Freeze calibration, model versions and rejection criteria before testing held-out shots, golfers or sessions appropriate to the deployment claim.
If a bias correction is fitted, report its uncertainty and validate it independently.
\item \textbf{Report effect sizes and uncertainty.}
Show paired differences, bias, spread, outliers and uncertainty intervals in physical units, stratified by relevant conditions.
Account for repeated shots within golfers and sessions when estimating uncertainty; hundreds of shots from one person are not hundreds of independent golfers.
Use predeclared tolerances tied to the coaching, fitting or research decision.
\end{enumerate}

\begin{implication}
Validation follows the whole chain from observation to decision.
Timing and geometry determine what is observed; the estimator determines what is inferred; the flight model determines what is predicted.
Log those distinctions with the uncertainty and version of every reported quantity, and preserve replayable observations where feasible.
That makes a disagreement investigable without assuming that the more expensive device, the smaller variance or the label ``measured'' identifies the correct answer.
\end{implication}
